\documentclass[11pt]{article}
\usepackage[a4paper,margin=2.5cm]{geometry}
\usepackage{CharisSIL}
\usepackage[hidelinks]{hyperref}
\setlength{\parindent}{0pt}
\setlength{\parskip}{0.75em}
\setlength{\emergencystretch}{2em}
\begin{document}
\thispagestyle{empty}

\textbf{21 September 2026}

Editor-in-Chief\\
\textit{Neurocomputing}

Dear Editor,

We submit the manuscript entitled \textbf{``Separating Clean Fidelity from Corruption Sensitivity in Quantized Neural Detectors: A Paired Evaluation''} for consideration as a research article in \textit{Neurocomputing}.

The paper addresses the evaluation of compressed neural networks after deployment. A corrupted-input accuracy gap can combine a clean-fidelity difference with a change associated with input degradation. Our executable paired four-arm protocol separates these quantities using controlled inputs, common image-bootstrap draws, and absolute task accuracy. Its contribution is a tested evaluation procedure and empirical characterization, rather than a new quantizer or a claim that one numerical format is universally superior.

Across six selection-disjoint VOC/KITTI blocks, FP8 has 1.60 AP points higher original-clean accuracy and 1.05 AP points higher corrupted accuracy than INT8, yet the clean-adjusted interaction is $-0.55$ AP (95\% paired-image percentile interval $[-0.78,-0.29]$). A four-dataset exploratory grid has 56 raw-versus-adjusted point-sign disagreements among 144 conditions. Controlled clean-input substitution, scale and covariance analyses, seed sensitivities, and recipe-portability tests delimit the result; they do not establish a population-wide format ranking. A retained prediction-level example reconstructs AP and uncertainty without cached aggregate scores.

The intended readership is researchers in neural network quantization, efficient inference, and evaluation of deployed learning systems. The manuscript relates its fixed-treatment evaluation question to recent \textit{Neurocomputing} work on Attention Round, stabilized activation-scale estimation, PLMQ, and QRT-DETR. Those approaches improve quantization recipes; our study examines how to interpret the accuracy and degradation sensitivity of the resulting executable candidates, using object detection as the measured setting.

An earlier version was declined by \textit{Computer Vision and Image Understanding} on readership/scope grounds; the editor suggested considering \textit{Neurocomputing}. This is a new submission for independent consideration, not an assurance of transfer or acceptance. The framing and related-work discussion have been revised accordingly.

Supplementary Information, highlights, and a graphical abstract accompany the manuscript. The experimental evidence is archived at \url{https://doi.org/10.5281/zenodo.22664869}; that release retains the preceding manuscript title and is not claimed to contain this revised text. AI assistance is disclosed, and the authors retain responsibility for the work.

Thank you for considering our work.

Sincerely,\\[0.8em]
\textbf{Anh Vu Le}\\
Corresponding author\\
Advanced Intelligent Technology Research Group\\
Faculty of Electrical and Electronics Engineering\\
Ton Duc Thang University, Ho Chi Minh City, Vietnam\\
\href{mailto:leanhvu@tdtu.edu.vn}{leanhvu@tdtu.edu.vn}

\end{document}
